\begin{theorem}[Cyclic-sector single-site grading]
    \label{thm:cyclic_sector_offdiag_shift}
    \lean{MPSTensor.IsCyclicSectorDecomp.offDiag_shift}
    \leanok
    \uses{def:periodic_tensor, def:cyclic_sector_decomp,
        thm:offdiag_shift_adjoint_cyclic_shift}
    If a periodic tensor admits a cyclic sector decomposition, then for every
    $k$ and every physical letter $i$ there exists a family
    $P_0,\ldots,P_{m-1}$ such that each $P_u$ is an orthogonal projection,
    \begin{align}
        P_u^\dagger&=P_u,\qquad P_u^2=P_u, \notag\\
        \sum_{u=0}^{m-1}P_u&=\Id, \notag
    \end{align}
    and $P_{k+1}A^i=A^iP_k$.
\end{theorem}

\begin{proof}\leanok
    Choose the projection family contained in the cyclic sector decomposition.
    It satisfies
    \begin{align}
        P_u^\dagger&=P_u,\qquad P_u^2=P_u, \notag\\
        \sum_uP_u&=\Id, \notag\\
        \E_A^\dagger(P_{u+1})&=P_u. \notag
    \end{align}
    Apply Theorem~\ref{thm:offdiag_shift_adjoint_cyclic_shift}.
\end{proof}

\begin{theorem}[Cyclic-sector off-diagonal reconstruction]
    \label{thm:cyclic_sector_eq_sum_offdiag}
    \lean{MPSTensor.IsCyclicSectorDecomp.eq_sum_offDiag}
    \leanok
    \uses{def:periodic_tensor, def:cyclic_sector_decomp,
        thm:eq_sum_offdiag_adjoint_cyclic_shift}
    If a periodic tensor admits a cyclic sector decomposition, then for every
    physical letter $i$ there exists a family $P_0,\ldots,P_{m-1}$ such that
    each $P_u$ is an orthogonal projection,
    \begin{align}
        P_u^\dagger&=P_u,\qquad P_u^2=P_u, \notag\\
        \sum_{u=0}^{m-1}P_u&=\Id, \notag
    \end{align}
    and
    \begin{align}
        A^i=\sum_{u=0}^{m-1}P_{u+1}A^iP_u. \notag
    \end{align}
\end{theorem}

\begin{proof}\leanok
    Choose the projection family contained in the cyclic sector decomposition.
    It satisfies
    \begin{align}
        P_u^\dagger&=P_u,\qquad P_u^2=P_u, \notag\\
        \sum_uP_u&=\Id, \notag\\
        \E_A^\dagger(P_{u+1})&=P_u. \notag
    \end{align}
    Apply Theorem~\ref{thm:eq_sum_offdiag_adjoint_cyclic_shift}.
\end{proof}

\begin{theorem}[Cyclic sector decomposition of a periodic tensor]
    \label{thm:cyclic_sector_decomp_after_blocking_periodic}
    \lean{MPSTensor.exists_cyclic_sector_decomp_after_blocking_of_isPeriodic}
    \leanok
    \uses{def:cyclic_sector_decomp, def:periodic_tensor}
    If $A$ is periodic with period $m$, then the blocked tensor $A^{(m)}$
    admits a cyclic sector decomposition. Thus there are projectors $P_u$ and
    corner tensors $C_u=P_u A^{(m)}$ which reproduce the blocked matrix-product
    vectors, are left-canonical, have nonzero bond dimensions, and realize the
    cyclic adjoint-transfer relation of Lemma~\emph{bdcf} in
    \cite[lines~404--423]{DeLasCuevas2017Irreducible}. The normality and
    non-repetition conclusions of Lemma~\emph{bdcf} are recorded separately in
    the component lemmas below.
\end{theorem}

\begin{proof}\leanok
    This is the cyclic-sector existence part of Lemma~\emph{bdcf}, translated
    into the inverse indexing convention of
    Definition~\ref{def:cyclic_sector_decomp}.
\end{proof}

\begin{theorem}[Cyclic-sector decomposition with compressed corner letters]
    \label{thm:cyclic_sector_compression_letter}
    \lean{MPSTensor.exists_cyclic_sector_decomp_with_letter_after_blocking_of_isPeriodic}
    \leanok
    \uses{thm:cyclic_sector_decomp_after_blocking_periodic, def:blocked_tensor}
    Under the hypotheses of the cyclic-sector construction for the blocked
    tensor $A^{(m)}$, there are sector tensors $C_k$, projections $P_k$, and
    linear corner isomorphisms
    \begin{align}
        \varphi_k\colon \MN{\dim C_k}
        \longrightarrow P_k\MN{D}P_k \notag
    \end{align}
    whose sector tensors are left-canonical, reproduce the blocked
    matrix-product vectors, have nonzero bond dimensions, and carry the cyclic
    adjoint-transfer decomposition of
    Theorem~\ref{thm:cyclic_sector_decomp_after_blocking_periodic}. In addition,
    for every blocked physical letter $i$,
    \begin{align}
        \varphi_k(C_k^i)=P_k(A^{(m)})^iP_k. \notag
    \end{align}
    This is the corner-letter identity used in
    \cite[Lemma~\emph{bdcf}, eq.~\emph{Cu}]{DeLasCuevas2017Irreducible}.
\end{theorem}

\begin{proof}\leanok
    Let $U$ diagonalize $P_k$, so that
    \begin{align}
        U^\dagger P_kU
        =
        \begin{pmatrix}
            I&0\\
            0&0
        \end{pmatrix}. \notag
    \end{align}
    Write $B^i=U^\dagger(A^{(m)})^iU$, and let $B_{kk}^i$ be the corner block
    on the support of $P_k$. The sector tensor is
    $C_k^i=V_k^\dagger B_{kk}^iV_k$, where $V_k$ identifies the support with
    $\C^{\dim C_k}$. The linear corner isomorphism is the inverse transport
    \begin{align}
        X\longmapsto
        U
        \begin{pmatrix}
            V_kXV_k^\dagger&0\\
            0&0
        \end{pmatrix}
        U^\dagger. \notag
    \end{align}
    Hence
    \begin{align}
        \varphi_k(C_k^i)
        &=
        U
        \begin{pmatrix}
            B_{kk}^i&0\\
            0&0
        \end{pmatrix}
        U^\dagger
        =P_k(A^{(m)})^iP_k. \notag
    \end{align}
\end{proof}

\begin{theorem}[Periodic cyclic sectors with support isometries]
    \label{thm:cyclic_sector_compression_letter_isometry_periodic}
    \lean{MPSTensor.exists_cyclic_sector_decomp_with_letter_and_isometry_after_blocking_of_isPeriodic}
    \leanok
    \uses{thm:cyclic_sector_decomp_after_blocking_isometry, def:periodic_tensor}
    If $A$ is periodic with period $m$, then the cyclic-sector decomposition of
    $A^{(m)}$ may be chosen with rectangular support isometries $V_k$ such that
    \begin{align}
        V_k^\dagger V_k&=\Id, \notag\\
        V_kV_k^\dagger&=P_k, \notag\\
        \varphi_k(X)&=V_kXV_k^\dagger. \notag
    \end{align}
    In particular, the same isometries realize the corner-letter identity
    $\varphi_k(C_k^i)=P_k(A^{(m)})^iP_k$.
    \begin{proof}
        \leanok
        Apply Theorem~\ref{thm:cyclic_sector_decomp_after_blocking_isometry}
        to the original tensor $A$ with its period-$m$ peripheral-spectrum
        hypotheses. The resulting decomposition is the decomposition of
        $A^{(m)}$, and the isometries $V_k$ in its conclusion are the required
        support isometries.
    \end{proof}
\end{theorem}

\begin{lemma}[One-site transport of sector overlaps]
    \label{lem:sector_overlap_succ_eq_cyclic_sector_decomp}
    \lean{MPSTensor.sectorOverlap_succ_eq_of_cyclicSectorDecomp}
    \leanok
    \uses{def:cyclic_sector_decomp, def:mpv_overlap, def:blocked_tensor,
        thm:offdiag_shift_adjoint_cyclic_shift}
    Let $A$ and $B$ be left-canonical tensors with cyclic sector
    decompositions $\{A_u\}_{u\in\Z_m}$ and
    $\{B_v\}_{v\in\Z_m}$. For every $u,v\in\Z_m$ and every $N>0$,
    $O_{A_{u+1},B_{v+1}}(N)=O_{A_u,B_v}(N)$.
\end{lemma}

\begin{proof}\leanok
    Choose the projection families $P_u$ and $Q_v$ from the two cyclic sector
    decompositions. The single-site relations
    $P_{u+1}A^i=A^iP_u$ and $Q_{v+1}B^i=B^iQ_v$ give the following
    trace-sum identity, where $T$ is the one-site rotation of physical words:
    \begin{align}
        O_{A_{u+1},B_{v+1}}(N)
        &=\sum_{\tau\in[d]^{Nm}}\tr(P_{u+1}A^\tau)
            \overline{\tr(Q_{v+1}B^\tau)} \notag\\
        &=\sum_{\tau\in[d]^{Nm}}\tr(P_{u+1}A^{T\tau})
            \overline{\tr(Q_{v+1}B^{T\tau})} \notag\\
        &=\sum_{\tau\in[d]^{Nm}}\tr(P_uA^\tau)
            \overline{\tr(Q_vB^\tau)}
        =O_{A_u,B_v}(N).
        \label{eq:pft_sector_overlap_transport}
    \end{align}
    The first and last equalities in
    (\ref{eq:pft_sector_overlap_transport}) are the cyclic-sector trace formula
    after reindexing blocked words as physical words of length $Nm$.
\end{proof}

\begin{lemma}[Off-diagonal eigenvectors of the period transfer vanish]
    \label{lem:offdiag_eigenvector_eq_zero_periodic}
    \lean{MPSTensor.offDiag_eigenvector_eq_zero_of_isPeriodic}
    \leanok
    \uses{def:periodic_tensor, def:cyclic_sector_decomp,
        lem:peripheral_pow_singleton}
    Let $A$ be periodic of period $m$. Let $P_0,\ldots,P_{m-1}$ be
    orthogonal projections satisfying
    \begin{align}
        \sum_{k=0}^{m-1}P_k&=\Id, \notag\\
        \E_A^\dagger(P_{k+1})&=P_k. \notag
    \end{align}
    If
    \begin{align}
        U&=P_uUP_v,\qquad P_uP_v=0, \notag\\
        \E_A^m(U)&=\zeta U,\qquad |\zeta|=1, \notag
    \end{align}
    then $U=0$.
\end{lemma}

\begin{proof}\leanok
    Since the peripheral spectrum of $\E_A$ consists of $m$-th roots
    of unity, Lemma~\ref{lem:peripheral_pow_singleton} gives
    $\spec_{\mathrm{per}}(\E_A^m)=\{1\}$, hence $\zeta=1$. The cyclic trace
    gives $\tr U=\tr(P_uUP_v)=\tr(P_vP_uU)=0$.

    Put $W=\sum_{t=0}^{m-1}\E_A^t(U)$. The equation
    $\E_A^m(U)=U$ gives $\E_A(W)=W$, and trace preservation gives
    $\tr W=0$. Irreducibility of $A$ therefore gives $W=0$. The shift relation
    $P_{k+1}A^i=A^iP_k$ moves $\E_A^t(U)$ into the $(u+t,v+t)$ corner, so only
    the $t=0$ term contributes after multiplying by $P_u$ and $P_v$. Thus
    $U=P_uWP_v=0$.
\end{proof}

% The overlap component lemmas below are part of the source dichotomy proof.
% The self-overlap and no-sector-match theorems are already formal statements;
% the full periodic overlap dichotomy still requires assembling these pieces
% with the remaining sector-match argument.
\begin{theorem}[Different periods imply asymptotic orthogonality]%
    \label{thm:periodic_overlap_zero_ne_period}
    \lean{MPSTensor.periodicOverlap_tendsto_zero_of_ne_period}
    \leanok
    \uses{def:periodic_tensor}
    If $A$ and $B$ have periods $m_A\ne m_B$, then
    $\lim_{N\to\infty}O_{AB}(N)=0$.
\end{theorem}
\begin{proof}\leanok
    \uses{thm:periodic_overlap_zero_ne_dim}
    If $D_A\ne D_B$, use Theorem~\ref{thm:periodic_overlap_zero_ne_dim}.
    Assume $D_A=D_B$ and $A\sim_{\mathrm{gp}}B$, say
    $B^i=\zeta XA^iX^{-1}$ with $X\in\GL_D(\C)$ and
    $\zeta\ne0$. The transfer maps satisfy
    \begin{align}
        \E_B(Y)
        =
        \zeta\overline{\zeta}
        X\E_A(X^{-1}YX^{-*})X^*. \notag
    \end{align}
    With $\E_A(\rho_A)=\rho_A$, put $\sigma=X\rho_AX^*$. Then
    $\E_B(\sigma)=|\zeta|^2\sigma$, while $\E_B(\rho_B)=\rho_B$ for a
    nonzero positive fixed point $\rho_B$. Perron--Frobenius eigenvalue
    uniqueness \cite[Theorem~6.3]{Wolf2012Quantum} gives $|\zeta|^2=1$.
    Hence $\E_B$ is conjugate to $\E_A$, so
    $\spec_{\mathrm{per}}(\E_A)=\spec_{\mathrm{per}}(\E_B)$. Since
    \begin{align}
        \spec_{\mathrm{per}}(\E_A)
        &=\{e^{2\pi ir/m_A}:0\le r<m_A\}, \notag\\
        \spec_{\mathrm{per}}(\E_B)
        &=\{e^{2\pi ir/m_B}:0\le r<m_B\},
        \label{eq:pft_periodic_set_equality}
    \end{align}
    equality of the two finite sets in (\ref{eq:pft_periodic_set_equality})
    gives $m_A=m_B$, a contradiction. Thus
    $A\not\sim_{\mathrm{gp}}B$, and the irreducible trace-preserving overlap
    dichotomy gives $\lim_{N\to\infty}O_{AB}(N)=0$.
\end{proof}

\begin{lemma}[Sector-block normality for periodic tensors]\label{lem:sector_blocked_normal_periodic}
    \lean{MPSTensor.sectorBlocked_isNormal_of_isPeriodic}
    \leanok
    \uses{def:periodic_tensor, def:cyclic_sector_decomp}
    For a periodic tensor equipped with a cyclic sector decomposition,
    each nonzero blocked sector tensor is normal.
\end{lemma}

\begin{theorem}[No sector match implies asymptotic orthogonality]%
    \label{thm:periodic_overlap_zero_no_sector_match}
    \lean{MPSTensor.periodicOverlap_tendsto_zero_of_no_sector_match}
    \leanok
    \uses{def:periodic_tensor, lem:sector_blocked_normal_periodic}
    Suppose $A$ and $B$ have the same total bond dimension and the same
    period $m$. Let $A_u$ and $B_v$ be period-blocked sector tensors of virtual
    dimensions $d_u$ and $e_v$, indexed by $u,v\in\Z_m$, such that
    \begin{align}
        \bar A&\sim\bigoplus_{u\in\Z_m}A_u, \notag\\
        \bar B&\sim\bigoplus_{v\in\Z_m}B_v, \notag
    \end{align}
    with unit weights, and the two decompositions are cyclic sector
    decompositions. Assume each sector is left-canonical:
    \begin{align}
        \sum_{\alpha}(A_u^\alpha)^\dagger A_u^\alpha&=I_{d_u}, \notag\\
        \sum_{\alpha}(B_v^\alpha)^\dagger B_v^\alpha&=I_{e_v}. \notag
    \end{align}
    Assume $d_u\ne0$ for every $u$ and $e_v\ne0$ for every $v$, and
    whenever $d_u=e_v$, the corresponding sector tensors are not gauge-phase
    equivalent after identifying their virtual spaces:
    \begin{align}
        d_u=e_v\quad\Longrightarrow\quad
        A_u\not\sim_{\mathrm{gp}}B_v. \notag
    \end{align}
    Then $\lim_{N\to\infty}O_{AB}(N)=0$.
\end{theorem}
\begin{proof}\leanok
    \uses{lem:sector_blocked_normal_periodic,
        thm:periodic_self_overlap_tendsto,
        thm:overlap_decay_irr_tp,
        thm:overlap_decay_rect_irr_tp}
    Write $\bar A\sim\bigoplus_{u\in\Z_m}A_u$ and
    $\bar B\sim\bigoplus_{v\in\Z_m}B_v$. For every $N$,
    \begin{align}
        O_{\bar A\bar B}(N)
        =
        \sum_{u\in\Z_m}\sum_{v\in\Z_m}O_{A_uB_v}(N).
        \label{eq:pft_overlap_double_sum}
    \end{align}
    For each sector pair $(u,v)$,
    \begin{align}
        d_u=e_v,\quad A_u\not\sim_{\mathrm{gp}}B_v
        &\quad\Longrightarrow\quad
        \lim_{N\to\infty}O_{A_uB_v}(N)=0, \notag\\
        d_u\ne e_v
        &\quad\Longrightarrow\quad
        \lim_{N\to\infty}O_{A_uB_v}(N)=0. \notag
    \end{align}
    The finite double sum (\ref{eq:pft_overlap_double_sum}) therefore gives
    \begin{align}
        \lim_{N\to\infty}O_{\bar A\bar B}(N)
        &=
        \sum_{u\in\Z_m}\sum_{v\in\Z_m}
        \lim_{N\to\infty}O_{A_uB_v}(N)
        =0. \notag
    \end{align}

    If $A\sim_{\mathrm{gp}}B$, then $\bar A\sim_{\mathrm{gp}}\bar B$, so for
    some $X\in\GL_D(\C)$ and $\zeta\ne0$,
    $(\bar B)^\alpha=\zeta X(\bar A)^\alpha X^{-1}$. Thus
    $V^{(N)}(\bar B)_\sigma=\zeta^N V^{(N)}(\bar A)_\sigma$, and
    \begin{align}
        O_{\bar A\bar B}(N)
        &=\overline{\zeta}^{\,N}O_{\bar A\bar A}(N), \notag\\
        O_{\bar B\bar B}(N)
        &=|\zeta|^{2N}O_{\bar A\bar A}(N). \notag
    \end{align}
    Since
    $\lim_{N\to\infty}O_{\bar A\bar A}(N)
      =\lim_{N\to\infty}O_{\bar B\bar B}(N)=m>0$,
    \begin{align}
        1
        &=
        \lim_{N\to\infty}
        \frac{O_{\bar B\bar B}(N)}{O_{\bar A\bar A}(N)}
        =
        \lim_{N\to\infty}|\zeta|^{2N}, \notag
    \end{align}
    hence $|\zeta|=1$. Therefore
    \begin{align}
        \lim_{N\to\infty}|O_{\bar A\bar B}(N)|
        =
        \lim_{N\to\infty}O_{\bar A\bar A}(N)
        =
        m\ne0, \notag
    \end{align}
    contradicting $\lim_{N\to\infty}O_{\bar A\bar B}(N)=0$. Thus
    $A\not\sim_{\mathrm{gp}}B$, and the irreducible trace-preserving overlap
    dichotomy gives $\lim_{N\to\infty}O_{AB}(N)=0$.
\end{proof}

\begin{lemma}[Sector-match propagation under translation]\label{lem:sector_match_propagation}
    \lean{MPSTensor.sectorMatch_propagation}
    \leanok
    \uses{def:periodic_tensor, def:cyclic_sector_decomp,
        lem:sector_overlap_succ_eq_cyclic_sector_decomp}
    Let $A$ and $B$ be left-canonical tensors with period $m$, and let
    $\bar A\sim\bigoplus_uA_u$ and $\bar B\sim\bigoplus_vB_v$ be unit-weight
    cyclic sector decompositions whose sectors are left-canonical. Suppose
    $d_{u_0}=e_{v_0}$, $d_{u_0}\ne0$, and, after identifying the virtual spaces,
    $A_{u_0}\sim_{\mathrm{gp}}B_{v_0}$. Then for every $l\in\Z_m$ there
    is an equality $d_{u_0+l}=e_{v_0+l}$ and, after this identification,
    $A_{u_0+l}\sim_{\mathrm{gp}}B_{v_0+l}$.
\end{lemma}

\begin{lemma}[Finite-cycle coboundary]\label{lem:finite_cycle_coboundary}
    \lean{TNLean.Algebra.exists_fin_succ_coboundary_of_prod_eq_one}
    \lean{TNLean.Algebra.exists_fin_cyclic_coboundary_of_prod_eq_one}
    \lean{TNLean.Algebra.exists_fin_complex_unit_cyclic_coboundary_of_prod_eq_one}
    \leanok
    \uses{}
    Let $G$ be a commutative group, and let
    $\kappa_0,\ldots,\kappa_n\in G$. If
    $\prod_{k=0}^{n}\kappa_k=1$, then there are
    $\varphi_0,\ldots,\varphi_n\in G$ such that
    $\kappa_k=\varphi_k\varphi_{k+1}^{-1}$ for every $k<n$, and
    $\kappa_n=\varphi_n\varphi_0^{-1}$.
    Equivalently, on a nonempty finite cyclic index set, a product-one family
    satisfies $\kappa_k=\varphi_k\varphi_{k+1}^{-1}$ for every cyclic edge.
    In particular, a product-one family of complex unit phases admits such a
    choice with all $\varphi_k$ again of unit modulus.
\end{lemma}

\begin{proof}\leanok
    \uses{}
    Set $\varphi_k=(\kappa_0\cdots\kappa_{k-1})^{-1}$, with the empty product
    equal to $1$. For $k<n$,
    \begin{align}
        \varphi_k\varphi_{k+1}^{-1}
        &=
        (\kappa_0\cdots\kappa_{k-1})^{-1}
        (\kappa_0\cdots\kappa_k)
        =\kappa_k. \notag
    \end{align}
    Separating the final factor in the product-one hypothesis gives
    $(\kappa_0\cdots\kappa_{n-1})\kappa_n=1$, which gives
    \begin{align}
        \varphi_n\varphi_0^{-1}
        =
        (\kappa_0\cdots\kappa_{n-1})^{-1}
        =
        \kappa_n. \notag
    \end{align}
\end{proof}

\begin{definition}[Global gauge]
    \label{def:periodic_global_gauge}
    \lean{MPSTensor.globalGauge}
    \leanok
    Given a positive period $m$, cyclic projectors $P_u$ and $Q_v$, a cyclic
    offset $q$, and corner operators $U_v$, the global gauge is
    \begin{align}
        U=\sum_uP_uU_{u+q}Q_{u+q}. \notag
    \end{align}
\end{definition}

\begin{lemma}[Global gauge as a sector sum]
    \label{lem:periodic_global_gauge_eq_sum}
    \lean{MPSTensor.globalGauge_eq_sum}
    \leanok
    \uses{def:periodic_global_gauge}
    Let $m\ge1$. If each corner operator satisfies
    $U_v=P_{v-q}U_vQ_v$, then $P_uU_{u+q}Q_{u+q}=U_{u+q}$, so the global
    gauge is the plain sector sum $U=\sum_vU_v$.
\end{lemma}

\begin{proof}\leanok
    \uses{def:periodic_global_gauge}
    The corner relation $U_v=P_{v-q}U_vQ_v$ gives
    $P_uU_{u+q}Q_{u+q}=U_{u+q}$; reindexing the sum by $v=u+q$ turns
    $\sum_uP_uU_{u+q}Q_{u+q}$ into $\sum_vU_v$.
\end{proof}

\begin{lemma}[Unit-phase gauge relation gives repeated blocks]
    \label{lem:gauge_phase_data_repeated_blocks}
    \lean{MPSTensor.repeatedBlocks_of_gaugePhaseData_norm_one}
    \leanok
    Suppose $B^i=\zeta XA^iX^{-1}$ and $|\zeta|=1$, with $X$ invertible.
    Then $A$ and $B$ satisfy the repeated-block relation
    $A^i=\zeta^{-1}X^{-1}B^iX$.
\end{lemma}
\begin{proof}\leanok
    Since $|\zeta|=1$, the scalar is nonzero. Multiplication by $X^{-1}$ on
    the left and by $X$ on the right gives
    $X^{-1}B^iX=\zeta A^i$; multiplication by $\zeta^{-1}$ proves the claim.
\end{proof}

\begin{lemma}[Global-gauge construction]
    \label{lem:periodic_global_gauge_assembly}
    \lean{MPSTensor.repeatedBlocks_of_globalGauge}
    \leanok
    \uses{def:cyclic_sector_decomp, def:periodic_global_gauge,
        lem:periodic_global_gauge_eq_sum,
        lem:left_canonical_irreducible_unitary_gauge}
    Let $m\ge1$, and let $A$ and $B$ carry the cyclic off-diagonal structure
    \begin{align}
        A^i&=\sum_uP_uA^iP_{u+1}, \notag\\
        B^i&=\sum_vQ_vB^iQ_{v+1}, \notag
    \end{align}
    with $P_u$ and $Q_v$ families of orthogonal projections summing to the
    identity. Let $q$ be the cyclic offset between the matched sector labels,
    and let each $U_v$ be a corner partial isometry from the $Q_v$ corner onto
    the $P_{v-q}$ corner, that is,
    \begin{align}
        U_v&=P_{v-q}U_vQ_v, \notag\\
        U_v^\dagger U_v&=Q_v, \notag\\
        U_vU_v^\dagger&=P_{v-q}. \notag
    \end{align}
    Assume the per-site proportionality of
    \cite[eq.~result]{DeLasCuevas2017Irreducible},
    \begin{align}
        P_uA^iP_{u+1}
        =
        \zeta\,U_{u+q}Q_{u+q}B^iQ_{u+q+1}U_{u+q+1}^\dagger,
        \qquad
        |\zeta|=1. \notag
    \end{align}
    Then the global gauge
    $U=\sum_uP_uU_{u+q}Q_{u+q}$ is unitary and satisfies
    $A^i=\zeta UB^iU^\dagger$ for every $i$, so $A$ and $B$ are repeated
    blocks.
\end{lemma}

\begin{proof}\leanok
    \uses{def:cyclic_sector_decomp, lem:periodic_global_gauge_eq_sum}
    The corner relation $U_v=P_{v-q}U_vQ_v$ gives
    $P_uU_{u+q}Q_{u+q}=U_{u+q}$, so $U=\sum_vU_v$ after reindexing. Distinct
    projections are orthogonal, $Q_vQ_w=0$ and $P_{v-q}P_{w-q}=0$ for
    $v\ne w$, which kills the cross terms in $UU^\dagger$ and $U^\dagger U$;
    the diagonal terms are $U_vU_v^\dagger=P_{v-q}$ and
    $U_v^\dagger U_v=Q_v$, so
    \begin{align}
        UU^\dagger&=\sum_uP_u=\Id, \notag\\
        U^\dagger U&=\sum_vQ_v=\Id. \notag
    \end{align}
    The off-diagonal structure of $B$ makes $Q_vB^iQ_w$ vanish unless
    $w=v+1$, so
    $UB^iU^\dagger=\sum_vU_vQ_vB^iQ_{v+1}U_{v+1}^\dagger$. Summing the
    per-site proportionality over $u$ and reindexing by $v=u+q$ gives
    \begin{align}
        \zeta UB^iU^\dagger
        &=
        \sum_v\zeta\,U_vQ_vB^iQ_{v+1}U_{v+1}^\dagger
        =
        \sum_uP_uA^iP_{u+1}
        =A^i. \notag
    \end{align}
    The last equality is the off-diagonal structure of $A$. The unit-modulus
    scalar and the unitary $U$ exhibit $A$ and $B$ as repeated blocks.
\end{proof}
